\documentclass[11pt]{article}

\usepackage[margin=1in]{geometry}
\usepackage{amsmath, amssymb, amsthm}
\usepackage{mathtools}
\usepackage{enumitem}

% ---- Semantics notation ----
\newcommand{\Exp}{\mathit{Exp}}
\newcommand{\Com}{\mathit{Com}}
\newcommand{\Store}{\mathit{Store}}
\newcommand{\conf}[2]{\langle #1, #2 \rangle}
\newcommand{\step}{\longrightarrow}
\newcommand{\steps}{\longrightarrow^{*}}
\newcommand{\nostep}{\not\longrightarrow}

\newcommand{\Case}[1]{\medskip\noindent\underline{\textbf{Case #1}}\smallskip\par}

% inference rule helper
\newcommand{\infrule}[2]{%
  \begin{array}{c}
    #1 \\ \hline #2
  \end{array}%
}
\newcommand{\infruleside}[3]{%
  \begin{array}{c}
    #1 \\ \hline #2
  \end{array}\quad #3%
}

\title{Homework 4}
\author{Alex Bastianini}
\date{}

\begin{document}
\maketitle

\section*{1)}
\begin{enumerate}
\item When $ c $ doesn't terminate or errors out.
\item If $ \models P \implies Q $.
\item When $ c $ terminates.
\end{enumerate}

\section*{2)}
\[
\infrule{ \{P \land b\} c \{Q\}\quad \{P \land \lnot b\} \text{ skip } \{Q\}}{ \{P\} \text{if } b \text{ then } c \{Q\}}
\]

We have proven during the exercise session that \texttt{if b then c} $ \sim $ \texttt{if b then c else skip}. Because the proposed rule is simply a special case of the Hoare rule applied to the equivalent command, it is sound. We could also further simplify the second condition to $ \models P \land \lnot b \implies Q $ using the \textsc{Skip} and \textsc{Consequence} rules.

\section*{3)}
\begin{enumerate}
  \item $ (z = m \land m \leq n) \lor (z = n \land n \leq m) $
  \item $ [(z = m - x \land m \leq n \land x \leq y) \lor (z = n - y \land n \leq m \land y \leq x)] \land x \geq 0 \land y \geq 0 $
\end{enumerate}

\section*{4)}
\[
\begin{array}{l}
\{ \mathrm{x} = i \land \mathrm{y} = j \land 0 \le i \} \implies \\
\{ \mathrm{x} \geq 0 \land \mathrm{y} = j \land 0 = (i-\mathrm{x}) \times j \}\\
\mathrm{z} := 0; \\
  \{ \mathrm{x} \geq 0 \land \mathrm{y} = j \land \mathrm{z} = (i-\mathrm{x}) \times j \}\\
\mathbf{while} \; (\mathrm{x} > 0) \; \mathbf{do} \; ( \\
  \{ \mathrm{x} \geq 0 \land \mathrm{y} = j \land \mathrm{z} = (i-\mathrm{x}) \times j \land \mathrm{x} > 0 \} \implies\\
  \{ \mathrm{x} - 1 \geq 0 \land \mathrm{y} = j \land \mathrm{z} + \mathrm{y} = (i-\mathrm{x} + 1) \times j \}\\
\quad \mathrm{z} := \mathrm{z} + \mathrm{y}; \\
  \{ \mathrm{x} - 1 \geq 0 \land \mathrm{y} = j \land \mathrm{z} = (i-\mathrm{x} + 1) \times j \}\\
\quad \mathrm{x} := \mathrm{x} - 1 \\
  \{ \mathrm{x} \geq 0 \land \mathrm{y} = j \land \mathrm{z} = (i-\mathrm{x}) \times j\}\\
); \\
  \{\mathrm{x} \geq 0 \land \mathrm{y} = j \land \mathrm{z} = (i-\mathrm{x}) \times j \land \lnot (\mathrm{x} > 0)\} \implies \\
\{ \mathrm{x} = 0 \land 0 = 0 \land \mathrm{z} = i \times j \}\\
\mathrm{y} := 0 \\
\{ \mathrm{x} = 0 \land \mathrm{y} = 0 \land \mathrm{z} = i \times j \}
\end{array}
\]
\section*{5)}
If we define $ I \triangleq (b \implies P) \land (\lnot b \implies Q) $, we can rewrite the rule using logical equivelences as follows
\[
  \infrule{\vdash \{I \land b\} c \{I\}}{\vdash \{I\} \text{while } b \text{ do } c \{I \land \lnot b\} }
\]
which is the classic \textsc{While} Hoare rule. Because Hoare rules are relatively complete, the equivalent \textsc{While-Alt} rule is too.

\end{document}
